\section{Domination and Douglas factorization}
\label{sec:domination}

\Cref{sec:direct_confinement} shows what follows once \(\AX\) vanishes.
To force the vanishing, the confinement theorem uses upper bounds on
\(\AX\) in the Loewner order.  In applications such bounds typically
come in two pieces: a main term, which an argument specific to the
application makes small, and an error term, kept explicit.

\begin{definition}[Completed domination bridge]
\label{def:duality_confinement:completed-domination-bridge}
Let \(\IOL\) be an involutive object ledger with
\(\int_X\|\psim\|^2\,d\mu<\infty\).  A \emph{completed domination
bridge} for \(\AX\) is a pair of positive trace-class operators
\(\Kminus\) and \(E\) on \(Y\) such that
\[
  \AX\preceq\Kminus+E.
\]
The operator \(\Kminus\) is the \emph{main term} and \(E\) the
\emph{defect}.  The bridge is \emph{exact} when \(E=0\).  A sequence of
bridges \(\AX\preceq\Kminus_n+E_n\) is a \emph{domination-record
sequence} if the bounds \(B_n=\Kminus_n+E_n\) satisfy
\(\trace B_n\to0\) as \(n\to\infty\).
\end{definition}

A bridge is bookkeeping rather than a theorem: it records which part of
an upper bound is the controlled main term and which part is error.
As a definition it imposes nothing beyond \(\AX\preceq\Kminus+E\);
indeed \(\Kminus=\AX\), \(E=0\) is always an exact bridge.  Its content
in applications is that \(\Kminus\) is computed from data that do not
involve \(\psim\).
The notation follows the framework's companion papers, where the main
term is computed on the carrier side from other data and the bridge
inequality is established by a separate argument.  For example, in the
finite-dimensional setting of \citep{TsiokosNeedleKiller2026},
\(\Kminus=L^-C^\dagger(L^-)^*\), where \(C\) is a positive audit
operator on a carrier space, \(C^\dagger\) its Moore--Penrose
pseudoinverse, and \(L^-\) an anti-invariant probe with
\(\ker C\subseteq\ker L^-\); the minus sign records that \(\Kminus\) is
the currency of the anti-invariant probe.  A
domination-record sequence is exactly the input of the confinement
theorem (\Cref{thm:duality_confinement:master-theorem}).  Exact
bridges have a useful reformulation, which rests on the following
classical theorem of Douglas.

\begin{theorem}[Douglas factorization {\citep{Douglas1966}}]
\label{thm:duality_confinement:douglas-domination}
Let \(H\), \(K_1\) and \(K_2\) be real Hilbert spaces and let
\(V:H\to K_1\) and \(W:H\to K_2\) be bounded linear operators.  Then
\[
  V^*V\preceq W^*W
  \quad\Longleftrightarrow\quad
  V=TW\ \text{for some bounded } T:K_2\to K_1 \text{ with } \|T\|\leq1.
\]
\end{theorem}

\begin{proof}
If \(V=TW\) with \(\|T\|\leq1\), then for every \(h\in H\),
\[
  \langle V^*Vh,h\rangle=\|TWh\|^2\leq\|Wh\|^2=\langle W^*Wh,h\rangle,
\]
so \(V^*V\preceq W^*W\).  Conversely, suppose \(V^*V\preceq W^*W\), so
that \(\|Vh\|\leq\|Wh\|\) for every \(h\in H\).  Define \(T_0\) on the
range of \(W\) by \(T_0(Wh)=Vh\).  It is well defined and linear: if
\(Wh=Wh'\) then \(\|Vh-Vh'\|\leq\|W(h-h')\|=0\).  It is contractive by
the same inequality.  Hence \(T_0\) extends by continuity to a
contraction on the closure \(\overline{\operatorname{Ran}W}\).  Setting
\(T=T_0\) on \(\overline{\operatorname{Ran}W}\) and \(T=0\) on its
orthogonal complement gives a bounded operator \(T:K_2\to K_1\) with
\(\|T\|\leq1\) and \(TW=V\).
\end{proof}

The theorem is due to Douglas \citep{Douglas1966}, who proves, for
bounded operators \(A,B\) on a single complex Hilbert space, that
\(AA^*\preceq\lambda^2BB^*\) for some \(\lambda\geq0\) if and only if
\(A=BC\) for a bounded operator \(C\), and that \(C\) can then be
chosen with \(\|C\|\leq\lambda\).  With \(\lambda=1\), \(A=V^*\) and
\(B=W^*\) this is the displayed statement after taking adjoints.  The
short proof above, which is standard, covers real scalars and
operators between different spaces.

\paragraph{Exact bridges as factorizations.}
By \Cref{rem:duality_confinement:analysis-operator},
\(\AX=V^*V\) for the analysis operator
\(V:Y\to L^2(X,\mu)\), \((Vh)(x)=\langle\psim(x),h\rangle\).  Every
positive operator \(\Kminus\) on \(Y\) can be written as
\(\Kminus=W^*W\), for instance with \(W=(\Kminus)^{1/2}\), which is
Hilbert--Schmidt when \(\Kminus\) is trace-class.
\Cref{thm:duality_confinement:douglas-domination} therefore says that
\(\AX\preceq\Kminus\) exactly when the analysis operator of the odd
readout factors as \(V=TW\) through a contraction \(T\).  In words: an
exact bridge exists precisely when every odd measurement
\(x\mapsto\langle\psim(x),h\rangle\) is obtained from the main-term
measurement \(Wh\) by one fixed contraction.
